% ═══════════════════════════════════════════════════════════════
% MUSTERLÖSUNG: El tiempo y los lugares
% Spanisch Klasse 8 - Unidad 6, DS 1
% ═══════════════════════════════════════════════════════════════

\documentclass[10pt, a4paper]{article}

\usepackage[utf8]{inputenc}
\usepackage[T1]{fontenc}
\usepackage[ngerman]{babel}

% --- SCHRIFTEN ---
\usepackage{palatino}
\usepackage{helvet}
\newcommand{\sanstext}[1]{{\fontfamily{phv}\selectfont #1}}

\usepackage{microtype}
\usepackage[margin=1.4cm, top=1.6cm, bottom=1.3cm]{geometry}
\usepackage{enumitem}
\usepackage{tabularx}
\usepackage{booktabs}
\usepackage{array}
\usepackage{multicol}
\usepackage{tikz}
\usetikzlibrary{calc}

\usepackage{fancyhdr}
\usepackage{lastpage}
\usepackage{needspace}

\usepackage{xcolor}
\usepackage{tcolorbox}
\tcbuselibrary{skins}

% ═══════════════════════════════════════════════════════════════
% FARBEN
% ═══════════════════════════════════════════════════════════════

\definecolor{kopfzeile}{HTML}{1A1A1A}
\definecolor{akzent}{HTML}{C41E3A}
\definecolor{hellgrau}{HTML}{F2F2F2}
\definecolor{mittelgrau}{HTML}{777777}
\definecolor{dunkelgrau}{HTML}{444444}
\definecolor{loesung}{HTML}{C62828}

% ═══════════════════════════════════════════════════════════════
% VARIABLEN
% ═══════════════════════════════════════════════════════════════

\newcommand{\thema}{El tiempo y los lugares}
\newcommand{\klasse}{8}
\newcommand{\maxpunkte}{28}

% ═══════════════════════════════════════════════════════════════
% KOPF- UND FUSSZEILE
% ═══════════════════════════════════════════════════════════════

\pagestyle{fancy}
\fancyhf{}
\renewcommand{\headrulewidth}{0pt}
\renewcommand{\footrulewidth}{0pt}
\cfoot{\sanstext{\footnotesize\textcolor{mittelgrau}{\thepage\,/\,\pageref{LastPage}}}}

% ═══════════════════════════════════════════════════════════════
% BOXEN
% ═══════════════════════════════════════════════════════════════

\newtcolorbox{vokabelkasten}{
    colback=white,
    colframe=mittelgrau,
    boxrule=0.5pt,
    arc=0pt,
    left=8pt, right=8pt, top=6pt, bottom=6pt
}

% ═══════════════════════════════════════════════════════════════
% BEFEHLE
% ═══════════════════════════════════════════════════════════════

\newcommand{\es}[1]{\textbf{#1}}
\newcommand{\de}[1]{{\small\textit{\textcolor{mittelgrau}{#1}}}}
\newcommand{\lsg}[1]{\textcolor{loesung}{\textbf{#1}}}
\newcommand{\punktebox}[1]{\hfill\sanstext{\small[\_\_/#1]}}

\newcommand{\aufgabe}[2]{%
    \needspace{4\baselineskip}%
    \vspace{12pt}%
    \noindent\sanstext{\textbf{#1}\quad#2}%
    \vspace{6pt}%
    \nopagebreak[4]%
}

\newlist{teil}{enumerate}{1}
\setlist[teil]{label=\alph*), leftmargin=1.8em, itemsep=4pt, parsep=0pt, topsep=3pt}

\setlength{\parindent}{0pt}
\setlength{\parskip}{4pt}
\setlength{\columnsep}{24pt}

% ═══════════════════════════════════════════════════════════════
% DOKUMENT
% ═══════════════════════════════════════════════════════════════

\begin{document}

% ═══════════════════════════════════════════════════════════════
% HEADER
% ═══════════════════════════════════════════════════════════════

\noindent
\begin{tikzpicture}[remember picture, overlay]
    \fill[kopfzeile] (current page.north west) rectangle +(\paperwidth, -1cm);
    \node[anchor=west, text=white, font=\sanstext\large\bfseries]
        at ([xshift=1.4cm, yshift=-0.5cm]current page.north west) {\thema{} -- \textcolor{loesung}{MUSTERLÖSUNG}};
    \node[anchor=east, text=white, font=\sanstext\small]
        at ([xshift=-1.4cm, yshift=-0.5cm]current page.north east)
        {Spanisch\,·\,Klasse \klasse\,·\,ML};
\end{tikzpicture}

\vspace{0.5cm}

% ═══════════════════════════════════════════════════════════════
% VOKABELKASTEN (mit Lösungen)
% ═══════════════════════════════════════════════════════════════

\begin{vokabelkasten}
\textbf{Kasten 1: El tiempo} \de{-- Das Wetter}

\vspace{6pt}

\begin{multicols}{2}
\begin{tabular}{@{}p{4cm}p{4cm}@{}}
¿Qué tiempo hace? & \lsg{Wie ist das Wetter?} \\[0.6em]
Hace sol. & \lsg{Es ist sonnig.} \\[0.6em]
Hace calor. & \lsg{Es ist heiß.} \\[0.6em]
Hace frío. & \lsg{Es ist kalt.} \\[0.6em]
Hace viento. & \lsg{Es ist windig.} \\
\end{tabular}

\columnbreak

\begin{tabular}{@{}p{4cm}p{4cm}@{}}
Hace buen tiempo. & \lsg{Das Wetter ist gut.} \\[0.6em]
Hace mal tiempo. & \lsg{Das Wetter ist schlecht.} \\[0.6em]
Llueve. & \lsg{Es regnet.} \\[0.6em]
Nieva. & \lsg{Es schneit.} \\[0.6em]
Hay nubes. & \lsg{Es ist bewölkt.} \\
\end{tabular}
\end{multicols}
\end{vokabelkasten}

\vspace{4pt}

\begin{vokabelkasten}
\textbf{Kasten 2: Los lugares} \de{-- Die Orte}

\vspace{6pt}

\begin{multicols}{3}
\begin{tabular}{@{}ll@{}}
el parque & \lsg{der Park} \\[0.5em]
la plaza & \lsg{der Platz} \\[0.5em]
el cine & \lsg{das Kino} \\[0.5em]
el teatro & \lsg{das Theater} \\
\end{tabular}

\columnbreak

\begin{tabular}{@{}ll@{}}
la biblioteca & \lsg{die Bibliothek} \\[0.5em]
la iglesia & \lsg{die Kirche} \\[0.5em]
el supermercado & \lsg{der Supermarkt} \\[0.5em]
la panadería & \lsg{die Bäckerei} \\
\end{tabular}

\columnbreak

\begin{tabular}{@{}ll@{}}
la farmacia & \lsg{die Apotheke} \\[0.5em]
el restaurante & \lsg{das Restaurant} \\[0.5em]
la cafetería & \lsg{das Café} \\[0.5em]
el centro comercial & \lsg{das EKZ} \\
\end{tabular}
\end{multicols}
\end{vokabelkasten}

\vspace{6pt}

% ═══════════════════════════════════════════════════════════════
% AUFGABE 1: Partner-Interview
% ═══════════════════════════════════════════════════════════════

\aufgabe{1}{Partner-Interview: ¿Qué haces cuando...?} \punktebox{12}

\vspace{4pt}

\begin{tabularx}{\textwidth}{|p{4cm}|X|}
\hline
\textbf{Frage} & \textbf{Mögliche Antworten (je 3P)} \\
\hline
\es{¿Qué haces cuando hace sol?} & \lsg{Cuando hace sol, voy al parque / a la playa / juego al fútbol con mis amigos.} \\[0.8em]
\hline
\es{¿Qué haces cuando llueve?} & \lsg{Cuando llueve, voy al cine / me quedo en casa / leo un libro.} \\[0.8em]
\hline
\es{¿Qué haces cuando hace frío?} & \lsg{Cuando hace frío, me quedo en casa / voy al centro comercial / tomo un chocolate caliente.} \\[0.8em]
\hline
\es{¿Qué haces cuando hace calor?} & \lsg{Cuando hace calor, voy a la piscina / tomo un helado / voy a la playa.} \\[0.8em]
\hline
\end{tabularx}

\vspace{4pt}
{\footnotesize\textit{Bewertung:} 3P pro vollständigem, korrektem Satz (1P Struktur, 1P Ort, 1P Grammatik)}

% ═══════════════════════════════════════════════════════════════
% AUFGABE 2: Wetter + Kleidung
% ═══════════════════════════════════════════════════════════════

\aufgabe{2}{¿Qué ropa llevas?} \punktebox{12}

\vspace{4pt}

\begin{teil}
    \item \es{¿Qué ropa llevas cuando hace mucho frío?}\\[0.3em]
    \lsg{Cuando hace mucho frío, llevo un abrigo, una bufanda, los guantes y el gorro.}

    \item \es{¿Qué ropa llevas cuando hace mucho calor?}\\[0.3em]
    \lsg{Cuando hace mucho calor, llevo pantalones cortos, una camiseta y las gafas de sol.}

    \item \es{¿Qué ropa llevas cuando llueve?}\\[0.3em]
    \lsg{Cuando llueve, llevo las botas y el paraguas / un impermeable.}

    \item \es{¿Qué ropa llevas cuando hace buen tiempo y vas al parque?}\\[0.3em]
    \lsg{Cuando hace buen tiempo, llevo una camiseta, vaqueros y zapatillas.}
\end{teil}

\vspace{4pt}
{\footnotesize\textit{Bewertung:} 3P pro Satz (1P Struktur cuando + llevo, 1P mind. 2 Kleidungsstücke, 1P korrekte Artikel)}

% ═══════════════════════════════════════════════════════════════
% AUFGABE 3: Beschreibe die Fotos
% ═══════════════════════════════════════════════════════════════

\aufgabe{3}{Beschreibe das Wetter auf den Fotos (Buch S. 102-103).} \punktebox{4}

\vspace{4pt}

{\footnotesize\textit{Hinweis: Die genauen Lösungen hängen von den Buchfotos ab. Beispielantworten:}}

\begin{multicols}{2}
\begin{tabular}{@{}ll@{}}
Foto 1: & \lsg{Hace sol / Hace buen tiempo.} \\[0.8em]
Foto 2: & \lsg{Llueve / Hace mal tiempo.} \\[0.8em]
Foto 3: & \lsg{Nieva / Hace mucho frío.} \\
\end{tabular}

\columnbreak

\begin{tabular}{@{}ll@{}}
Foto 4: & \lsg{Hace calor / Hace sol.} \\[0.8em]
Foto 5: & \lsg{Hay nubes / Hace viento.} \\[0.8em]
Foto 6: & \lsg{Hace frío / Nieva.} \\
\end{tabular}
\end{multicols}

\vspace{4pt}
{\footnotesize\textit{Bewertung:} Je 0,5-1P pro korrekter Wetterbeschreibung (flexibel nach Foto)}

% ═══════════════════════════════════════════════════════════════
% FOOTER
% ═══════════════════════════════════════════════════════════════

\vfill

\noindent\rule{\textwidth}{0.5pt}

\vspace{4pt}

\hfill\sanstext{\textbf{Maximalpunktzahl:}} \maxpunkte

\end{document}
